First optimize the separated-pilot envelope. Take power tunings
\[
h\asymp n^{-\gamma},
\quad b_\mu\asymp n^{-r_\mu},
\quad \delta_\mu\asymp n^{-q_\mu},
\quad b_\pi\asymp n^{-r_\pi},
\quad \delta_\pi\asymp n^{-q_\pi}.
\]
For the regression pilot, the exponent of \(u+h^{A_0}\) in Lemma~\texttt{lem:localized-polynomial-master-risk} cannot exceed
\[
\min\left\{S_0q_\mu,A_0r_\mu,
\frac{1-r_\mu-dq_\mu}{2},A_0\gamma\right\}.
\]
If the first three entries are all at least \(e\), then
\[
r_\mu\ge \frac{e}{A_0},
\qquad q_\mu\ge \frac{e}{S_0},
\qquad 1-r_\mu-dq_\mu\ge 2e.
\]
Adding the resulting lower bounds on \(r_\mu\) and \(q_\mu\) to the last inequality gives
\[\ne\left(2+A_0^{-1}+dS_0^{-1}\right)\le 1.
\]
Therefore the largest possible regression-pilot exponent at fixed \(\gamma\) is at most
\[
e_\mu:=\min\{A_0\gamma,E_\mu\}.
\]
This value is attained by taking
\[
r_\mu=\frac{e_\mu}{A_0},
\qquad q_\mu=\frac{e_\mu}{S_0},
\]
because the definition of \(E_\mu\) implies
\[
\frac{1-r_\mu-dq_\mu}{2}\ge e_\mu.
\]
The same calculation for the density pilot shows that its largest attainable exponent at fixed \(\gamma\) is
\[
e_\pi:=\min\{B_0\gamma,E_\pi\},
\]
attained at \(r_\pi=e_\pi/B_0\) and \(q_\pi=e_\pi/S_0\). For every \(0<\gamma<1\), these choices have strictly positive count exponents. Consequently all tuning restrictions in Lemma~\texttt{lem:localized-polynomial-master-risk} hold for all sufficiently large \(n\). Its exponential term is super-polynomial because \(nh\asymp n^{1-\gamma}\), and substitution gives the polynomial exponent
\[
\min\left\{2\alpha\gamma,1-\gamma,2(e_\mu+e_\pi),1\right\}.
\]
The entry \(1\) is redundant because \(1-\gamma\le1\). The remaining objective is continuous on \([0,1]\), is positive at some interior point, and is zero at both endpoints. Hence its maximum is achieved at an admissible interior value of \(\gamma\). The preceding upper-envelope argument proves that no other power tuning of this master-risk bound yields a larger exponent. This establishes the displayed formula for \(R_{\mathrm{sep}}\).

Next optimize the regressogram bound. Lemma~\texttt{lem:explicit-histogram-window-upper} gives, for \(h\asymp n^{-u}\) and \(\delta\asymp n^{-v}\),
\[
R(u,v)=\min\{2A_0u,4S_0v,2S_0v+2B_0u,1-u-dv,1\}.
\]
Fix a candidate exponent \(r>0\) and write \(x=u/r\) and \(y=v/r\). The first three inequalities required for \(R(u,v)\ge r\) are
\[
x\ge\frac{1}{2A_0},
\qquad y\ge\frac{1}{4S_0},
\qquad B_0x+S_0y\ge\frac12.
\]
At the lower corner \(x_0=(2A_0)^{-1}\), \(y_0=(4S_0)^{-1}\), the deficit in the third inequality is
\[
\Delta_*=\left(\frac12-B_0x_0-S_0y_0\right)_+
=\frac12\left(\frac12-\frac{B_0}{A_0}\right)_+.
\]
Filling one unit of this deficit by increasing \(x\) costs \(1/B_0\) in the objective \(x+dy\), whereas filling it by increasing \(y\) costs \(d/S_0\). Since this is a linear program with one residual constraint, its minimum normalized bandwidth cost is
\[
C_*=\frac1{2A_0}+\frac d{4S_0}+\left(\frac12-\frac{B_0}{A_0}\right)_+
\min\left\{\frac1{2B_0},\frac d{2S_0}\right\}.
\]
The variance inequality is \(u+dv\le1-r\), equivalently \(rC_*\le1-r\). Thus the largest feasible exponent is
\[
r=(1+C_*)^{-1}=R_{\mathrm{hist}}.
\]
At an optimizer, \(u+dv=1-R_{\mathrm{hist}}<1\), so the count restriction in Lemma~\texttt{lem:explicit-histogram-window-upper} holds for every sufficiently large \(n\). This also proves the asserted maximality of \(R_{\mathrm{hist}}\) within that power-tuned bound.

For each fixed admissible model tuple, use the separated-pilot estimator or the regressogram according as \(R_{\mathrm{sep}}\ge R_{\mathrm{hist}}\) or \(R_{\mathrm{hist}}>R_{\mathrm{sep}}\). Definition~\texttt{def:minimax-mse} and the two preceding risk bounds then give constants \(C_{\mathrm{up}}<\infty\) and \(n_0\), independent of \(n\), such that
\[
\mathcal M_n(\alpha,\beta,s,d,M_Y,L,c,t_0)
\le C_{\mathrm{up}}n^{-R_{\mathrm{ach}}}
\]
for all \(n\ge n_0\). Lemma~\texttt{lem:all-beta-oracle-lower-floor} and Definition~\texttt{def:oracle-rate} give a constant \(C_{\mathrm{low}}>0\), independent of \(n\), for which
\[
\mathcal M_n(\alpha,\beta,s,d,M_Y,L,c,t_0)
\ge C_{\mathrm{low}}n^{-2\alpha/(2\alpha+1)}.
\]
After enlarging \(n_0\) if necessary, these are the claimed two sides of the bracket.

It remains to characterize when the separated-pilot exponent reaches the oracle exponent. For every \(\gamma\in[0,1]\),
\[
\min\{2\alpha\gamma,1-\gamma\}
\le\frac{2\alpha}{2\alpha+1},
\]
and equality holds only at
\[
\gamma_0=(2\alpha+1)^{-1}.
\]
It follows that \(R_{\mathrm{sep}}=2\alpha/(2\alpha+1)\) if and only if the third entry of its objective is at least the oracle exponent at \(\gamma_0\), namely
\[
\min\{A_0\gamma_0,E_\mu\}+\min\{B_0\gamma_0,E_\pi\}
\ge\alpha\gamma_0.
\]
Elementary interval feasibility shows that this inequality holds if and only if there exists
\[
z\in[(\alpha-B_0)_+,A_0]
\]
such that
\[
E_\mu\ge z\gamma_0,
\qquad E_\pi\ge(\alpha-z)\gamma_0.
\]
Indeed, the first minimum can allocate at most \(A_0\gamma_0\) and the second at most \(B_0\gamma_0\), while their respective additional caps are \(E_\mu\) and \(E_\pi\).

Put \(C_0=2\alpha+1\). Direct substitution of the definitions of \(E_\mu\) and \(E_\pi\) shows that the first cap inequality is equivalent to
\[
D_\mu(z):=C_0-z/A_0-2z>0,
\qquad
S_0\ge\frac{dz}{D_\mu(z)},
\]
and the second is equivalent to
\[
D_\pi(z):=C_0-(\alpha-z)/B_0-2(\alpha-z)>0,
\qquad
S_0\ge\frac{d(\alpha-z)}{D_\pi(z)}.
\]
The strict positivity assertions also cover the endpoint cases: if \(z=0\), then \(D_\mu(z)=C_0>0\), and if \(\alpha-z=0\), then \(D_\pi(z)=C_0>0\); otherwise positivity follows because the right-hand side before division is strictly positive. By Definition~\texttt{def:oracle-threshold} and Lemma~\texttt{lem:separated-threshold-optimization}, such a \(z\) exists exactly when
\[
A_0+B_0\ge\alpha,
\qquad S_0\ge g(\alpha,\beta,d).
\]
Definition~\texttt{def:oracle-region} identifies these two conditions with \((\alpha,\beta,s,d)\in\Omega\). This proves the claimed equivalence.

Finally, because \(d/(4S_0)>0\), one has
\[
R_{\mathrm{hist}}<\left(1+\frac1{2A_0}\right)^{-1}.
\]
If \(\alpha\le1\), then \(A_0=\alpha\), and the right-hand side equals \(2\alpha/(2\alpha+1)\). If \(\alpha>1\), then \(A_0=1\), so
\[
R_{\mathrm{hist}}<\frac23<\frac{2\alpha}{2\alpha+1}.
\]
Thus \(R_{\mathrm{hist}}\) is always strictly below the oracle exponent. On the complement of \(\Omega\), the proved equivalence also makes \(R_{\mathrm{sep}}\) strictly smaller than the oracle exponent. By Definition~\texttt{def:oracle-region}, that complement is precisely the union of the two residual regions displayed in the statement, so their maximum \(R_{\mathrm{ach}}\) is strictly smaller there as well.

If \(B_0\ge A_0/2\), then the positive-part term in \(R_{\mathrm{hist}}\) vanishes, leaving no occurrence of \(\beta\). If \(\beta\ge1\), then \(B_0=1\), so \(E_\pi\), \(R_{\mathrm{sep}}\), and \(R_{\mathrm{hist}}\), and hence \(R_{\mathrm{ach}}\), are all constant as functions of the numerical value of \(\beta\). These are properties of the displayed achievable exponents only and do not supply a residual-region minimax converse.